% 05_physical.tex - the soil zone: parameters, the :SoilZone command, the run, the summary and the warnings
\section{Let the water table rise into the soil}

\step{Add soil parameters to \texttt{Creek.rvp}}
The soil zone will be \cmd{TOPSOIL} and \cmd{SUBSOIL} (\cmd{SOIL[0]} and \cmd{SOIL[1]}). The coupling needs three things:
\begin{itemize}
\item Recharge leaves from the bottom layer of the soil zone, \cmd{SUBSOIL}, by a vertical process. This model already
  has one: \cmd{:Percolation \dots SUBSOIL GROUNDWATER}.
\item Each soil-zone layer needs \cmd{POROSITY} and \cmd{FIELD\_CAPACITY}. The model has them.
\item The class of the bottom layer needs \cmd{HYDRAUL\_COND} [mm/d], the hydraulic conductivity of the soil. With
  \cmd{K\_VERT} of the top aquifer layer, it sets how easily water moves between soil and aquifer. Capillary rise (water
  drawn up from the water table into the soil) is on by default with unsaturated columns (Section~\ref{sec:vz}) and uses
  \cmd{CLAPP\_B} and \cmd{AIR\_ENTRY\_PRESSURE} [mm], two parameters of the soil's water retention curve.
\end{itemize}
Add the three columns to the soil parameter table of \cmd{Creek.rvp}:
\diff{t3_soil.diff}

\step{Switch the soil zone on in \texttt{Creek.rvg}}
Add one command at the end of \cmd{Creek.rvg}:
\diff{t3_physical.diff}

\step{Run Raven}\label{st:physrun}
Run into a new output folder:
\term{t3_run_physical.txt}
\see \msg{Exiting Gracefully: Successful Simulation}. The run takes about \V{t3:secphysical}~seconds instead of
\V{t3:secconceptual}, because every day the soil zone and MODFLOW are computed again and again until they agree. The files
of this step are in \cmd{\$WS/model/step\_06}.
\begin{trouble}
\begin{itemize}
\item \emph{``soil class \dots of SOIL[1] \dots needs HYDRAUL\_COND''}: the change to \cmd{Creek.rvp} is missing.
\item \emph{``RECHARGE\_DELAY \dots cannot be used with :SoilZone PHYSICAL''}: a profile has a recharge delay. Set it to 0;
  use \cmd{:SoilZoneVadose LAG} for a delay instead.
\item \emph{``soil-zone interface did not converge''} (rare): raise \cmd{:SoilZoneIteration}. If that does not
  help, increase \cmd{SEEPAGE\_SMOOTHING\_DEPTH} or use a shorter time step.
\item Compare your files with those of this step. This command prints the lines that differ, leaving spacing and
  empty lines aside:
\begin{shell}
$ diff -wB $WS/model/step_06/Creek.rvg Creek.rvg
\end{shell}
\end{itemize}
\end{trouble}

\step{Read the summary and the warnings}
\term{t3_summary.txt}
\why \V{t3:nszhru} HRUs have a soil zone. The lake HRU has no soil and keeps the seepage drain. The soil zone is 1.3~m
thick under the valley and the wetland and 1.0~m under the slopes (from the soil profiles of the \cmd{.rvp}). Raven keeps
one soil-zone column for every link between an HRU and a cell. The run wrote \V{t3:nwarn} warnings. Read them:
\term{t3_warnings.txt}
\begin{itemize}
\item \emph{The soil-zone base lies up to \V{t3:offmax}~m from the model top on \V{t3:offlinks} links.} The top of a cell
  is the land surface minus the soil thickness of the HRUs in it, weighted by their overlap with the cell. A cell shared by
  the valley (1.3~m) and a slope (1.0~m), or by an HRU with soil and the lake (no soil), has a top between the two. The
  warning says where. With HRUs whose soil profiles differ, this is expected.
\item \emph{The specific yield differs from the drainable porosity of the bottom soil layer on \V{t3:sylinks} links}
  (ratio down to \V{t3:syratio}). The till of the slopes has $S_y=0.08$. The subsoil has a drainable porosity of
  $0.40\times(1-0.25)=0.30$. A water table rising from the till into the soil would change speed there. On these slopes
  the water table stays far below the soil, so this does not matter here. On wet ground, make the two values agree.
\end{itemize}
